\documentclass[11pt,a4paper]{article}
\usepackage{amsmath,amssymb,amsthm,hyperref}
\usepackage[margin=27mm]{geometry}
\newtheorem{proposition}{Proposition}
\newcommand{\Q}{\mathbb Q}
\title{Prime-pattern scope for one- and two-square height covers}
\author{}\date{17 November 2026}
\begin{document}\maketitle

\section*{The one-square restriction}
Sun's Theorem 4.1 in arXiv:2610.20691v3 supplies patterns
$u+c_jw^2$ with a free linear variable $u$. Proposition 5.1 proves
the corresponding weighted average over $1\le u\le N$ and
$1\le w\le T$, with an error $o(NT)$. The change of variables
in its equation (5.13) retains a free translated coordinate $x$ in
$x+R_j(m)$. These statements do not impose $u-Aw^2$ to be a square.

In particular the forms $m^2+(A-\theta_j)n^2$ required for a
one-square height cover are not a specialization justified by those
theorems. In the relevant box, a restriction
$u=m^2+Aw^2$ selects at most $O(\sqrt N T)$ pairs from the
$NT$ pairs in the original average, for fixed $A$ and
$T=o(\sqrt N)$. The quoted error can exceed the entire restricted
count. Positivity of the unrestricted average therefore gives no
existence statement on this thin set. A separate restricted prime-pattern
theorem, including the conjugate reduction-order conditions, is needed.

This is a statement of the cited theorem's scope; it is not a proof
that the desired restricted patterns fail or that no stronger theorem
can establish them.

\section*{A two-square preparation that is within the theorem}
Let $M>1$ be a fixed positive rational, and retain the six contacts
$\Theta=\{0,-1,-1/2,-1/3,-1/4,-2/5\}$. The following assertion
uses precisely the retained prime-pattern theorem.

\begin{proposition}
The source prime-pattern preparation can impose simultaneously
\[
                 t-M=z_1^2+z_2^2,\qquad z_1z_2\ne0,
\]
with rational $t,z_1,z_2$, while retaining the six prime contacts,
their conjugate reduction-order compatibility, arbitrary fixed
denominator divisibility, finite avoidance, and a prescribed lower
bound for $t$.
\end{proposition}
\begin{proof}
Choose a positive integer $D$ with $D^2\theta$ and $D^2M$ integral
for all $\theta\in\Theta$. Apply Sun's Theorem 4.1 to the seven
distinct coefficients $-D^2\theta$ and $-D^2M$, including the
same fixed avoidance modulus containing three and seventeen.
Its lower bound for $u/w^2$ may exceed $D^2M$ and any other
fixed bound. Put $t=u/(D^2w^2)$.
The additional value
\[
                     p_*=u-D^2Mw^2
\]
is a positive prime congruent to one modulo eight. Its Gaussian
prime generator writes $p_*=a^2+b^2$ with nonzero integers $a,b$:
a zero component would make the prime a square. Therefore
\[
                  z_1=a/(Dw),\qquad z_2=b/(Dw)
\]
give the required identity. The same theorem includes all the
requested order and preparation conditions; adding one fixed
coefficient does not change its quantified scope. One may use only
the original six primes when selecting the shared auxiliary index.
\end{proof}

\section*{Guard and coordinate audit; unresolved height step}
For the proposed single-detector construction take
$M=(\beta^e+T^e)^{dm}$ after fixing the actual root point $A=[n]P$.
This scalar depends only on the retained rational abscissas; it has
the same value at every point-algebra character. Retain $t,z_1$ as
free coordinates and adjoin the finite generator
\[
                    Z_2^2=t-M-z_1^2.
\]
At the constructed tuple its radicand is the same nonzero rational
square at every character. No extra inverse witness is needed.
The intrinsic two-parameter pole block adds one further finite
quadratic generator. The rank is consequently $64\cdot2\cdot2=256$,
and the unrestricted-coordinate census is
\[
  1\text{ input}+6\text{ point coordinates}+1\text{ argument}
  +1\text{ free square coordinate}+2\text{ pole parameters}
  +1\text{ encoder}=12.
\]

This is not a twelve-variable definition. The existing mixed-power
height inequality applies to $t=a^d+z^2$ and has not been proved
for $t=a^d+z_1^2+z_2^2$. The real inequality $t\ge a^d$ does
not control the denominator contribution to $h(a)$. Thus the
prime-pattern and guard construction does not close the soundness
argument. A separate valid height theorem would be necessary.

\begin{thebibliography}{9}
\bibitem{Sun}Z.-W. Sun, \emph{48 unknowns for Hilbert's tenth problem
over the rationals}, arXiv:2610.20691v3 (11 November 2026),
Theorem 4.1, Proposition 5.1, equation (5.13).
\url{https://arxiv.org/pdf/2610.20691}.
\bibitem{TZ}T. Tao and T. Ziegler, \emph{Polynomial patterns in the
primes}, Forum of Mathematics, Pi 6 (2018), e1.
\url{https://doi.org/10.1017/fmp.2017.3}.
\end{thebibliography}
\end{document}
